% Generated draft: requires mathematical and source review.
\documentclass[11pt]{article}
\usepackage[margin=25mm]{geometry}
\usepackage{fontspec}
\setmainfont{Latin Modern Roman}
\usepackage{amsmath,amssymb,mathtools,unicode-math}
\setmathfont{Latin Modern Math}
\usepackage{enumitem,needspace,xurl,hyperref}
\setlength{\parindent}{0pt}
\setlength{\parskip}{5pt}
\setlength{\emergencystretch}{3em}
\newcommand{\N}{\mathbb{N}}
\newcommand{\Z}{\mathbb{Z}}
\newcommand{\Q}{\mathbb{Q}}
\newcommand{\R}{\mathbb{R}}
\newcommand{\C}{\mathbb{C}}
\newcommand{\F}{\mathbb{F}}
\newcommand{\A}{\mathbb{A}}
\newcommand{\PP}{\mathbb{P}}
\newcommand{\E}{\mathbb{E}}
\newcommand{\eps}{\varepsilon}
\newcommand{\dd}{\,\mathrm{d}}
\begin{document}
\section*{UnsolvedMath ID: 30002478}
\textbf{Source problem code:} catalogue entry
\subsection{Definitions and mathematical statement}
Under the usual complex-scalar convention, a \(C^*\)-algebra \(A\) is a complex Banach algebra with a conjugate-linear involution \(a\mapsto a^*\) satisfying \((ab)^*=b^*a^*\) and \(\|a^*a\|=\|a\|^2\). It is simple if it is nonzero and its only norm-closed two-sided ideals are \(\{0\}\) and \(A\). Infinite dimensionality means \(\dim_{\C}A=\infty\). Let \(\operatorname{Aut}(A)\) be the group of all complex-linear \(*\)-automorphisms of \(A\). Define its conjugacy equivalence relation by
\[
\alpha\sim_A\beta
\quad\Longleftrightarrow\quad
\exists\gamma\in\operatorname{Aut}(A)\;
\beta=\gamma\circ\alpha\circ\gamma^{-1},
\qquad \alpha,\beta\in\operatorname{Aut}(A).
\]
The supplied question is whether
\[
\exists\,A\quad
\bigl[A\text{ is an infinite-dimensional simple }C^*\text{-algebra and }
\sim_A\text{ is classifiable}\bigr].
\]
Here ``classifiable'' is retained verbatim as an unresolved qualification: the supplied record specifies neither the permitted invariants nor the regularity required of a classification map. Thus the display is a question schema, not a fully specified mathematical assertion. In particular, the record does not distinguish smooth classification from classification by countable structures, or another classification notion. It also imposes no explicit separability or unitality hypothesis. Such hypotheses must not be inserted without source clarification; if a descriptive-set-theoretic interpretation is intended, the relevant topology and measurable structures must also be specified.
\subsection{Short English statement}
Does some infinite-dimensional simple \(C^*\)-algebra have a classifiable conjugacy relation on its automorphism group? The intended meaning of ``classifiable'' and any additional hypotheses on the algebra remain unspecified in the supplied record.
\subsection{Sources}
\begin{itemize}
\item Supplied local UnsolvedMath registry record, permanent ID \(30002478\), source problem code \texttt{catalogue entry}, in the Oberwolfach Reports Open Problems collection. The record gives the proposed year as 2014 and supplies no named proposer.
\item Supplied report citation: \emph{Set Theory} (2014), Oberwolfach Reports, pp.~91--144, DOI \texttt{10.4171/owr/2014/02}. \url{https://doi.org/10.4171/owr/2014/02}. The supplied background describes Oberwolfach Reports as documenting mathematical workshops and their open-problem discussions; no exact problem page is supplied.
\item Supplied literature source, ``Oberwolfach set theory report'': \url{https://ems.press/content/serial-article-files/46488}.
\item Supplied literature source, ``Lupini, classification of C-star automorphisms'': \url{https://arxiv.org/abs/1501.03022}. This label and link are reproduced as supplied; their bibliographic correspondence has not been independently verified.
\item Literature assessment dated exactly \texttt{2026-08-21}: the supplied curation labels the problem open and states that no general classification of conjugacy in automorphism groups of infinite-dimensional simple \(C^*\)-algebras is known, while existing work gives substantial nonclassifiability results without settling the stated existence question. This is supplied dated curation derived from the OWR record, not an independently verified assessment. The supplied background explicitly requests independent verification of time-sensitive status and cited claims.
\end{itemize}
\end{document}
